%Derivative and fit

Furthermore, to determine the magnetic susceptibility $\chi_m$ of the test sample. We use Equations.~(\ref{eq:translationalForce}) and (\ref{eq:deflectionAngle}), 

\begin{subequations}
	\begin{equation}
		\chi_m = \frac{\mu_0 m g}{V}\frac{\tan (\alpha)}{ B_0 \frac{dB_0}{dz}},
		\label{eq:chi}
	\end{equation}
	\begin{equation}
		\frac{\chi_m}{\rho} = \mu_0 g\frac{\tan (\alpha)}{ B_0 \frac{dB_0}{dz}}.
		\label{eq:chi_rho}
	\end{equation}
	\label{eq:chis}
\end{subequations}

\noindent providing a direct estimation of magnetic susceptibility from angular deflection measurements, provided the field and its gradient are known. For experimental purposes, a plot of the deflection angle versus the product of field strength and field gradient allows one to extract $\chi_m$ from the slope of a linear fit, following

\begin{equation}
	\tan (\alpha) =   \frac{\chi_m }{\rho \mu_0 g}  \cdot B_0 \frac{dB_0}{dz}.
	\label{eq:forFit}
\end{equation}

In general, these forces are described as a function of the static field strength ($B_0$), its spatial gradient ($\nabla B_0$), and the material properties such as magnetic susceptibility ($\chi$) and density ($\rho$). This equation later allows for plotting and interpretation.